\section{Finite component bounds and constructive skeletons}
\label{sec:groupoid-skeletons}

Finite category tables have an explicit list of all their objects. A
primitively recursively presented category may instead have natural
numbers as objects, even when a finite numerical bound on its
isomorphism classes is known. Extracting representatives in this setting
is a genuine mathematical operation. The distinction also appears in
the inspected Mordell--Weil packet: its PR rational-halving routine
decides individual instances of division by two, while its generator
theorem requires a supplied complete list of representatives modulo
doubling~\cite{repo-mw}.

\subsection{A uniformly decidable two-class obstruction}

Fix a standard finite coding of Turing machines, configurations and
transitions. Let $h(e,n)\in\{0,1\}$ be one exactly when the machine
with code $e$, on blank input, has halted within $n$ steps. This is a
fixed primitive recursive function: simulate exactly $n$ transitions,
freezing a halted configuration. Define
\[
                 m\,E_e\,n\quad\Longleftrightarrow\quad
                            h(e,m)=h(e,n).
\]
For every $e$, this is a decidable equivalence relation. Its reflexivity,
symmetry and transitivity are equality laws for Boolean values. Among
any three natural numbers, two have the same Boolean value, so no three
pairwise inequivalent objects exist.
This is a counting bound, expressed by the failure of any three-element
list to be pairwise inequivalent; no covering list is included in it.

Define a thin groupoid $\mathcal G_e$ with objects $\mathbb N$ and an
arrow $m\longrightarrow n$ precisely when $E_e(m,n)$. Such an arrow
is encoded by its endpoint pair. Its identity is $(m,m)$, its inverse
is $(n,m)$, and the composite of $(m,n)$ followed by $(n,p)$ is
$(m,p)$. Validity is checked by $E_e$, and the groupoid laws reduce
to equality of endpoint pairs. Hence this is a uniform PR family of
presentations, with at most two isomorphism classes. It uses only one
fixed bounded machine simulator, not a universal total evaluator for
primitive recursive functions.

\begin{theorem}[No uniform finite representative producer]
\label{thm:no-two-class-skeleton}
There is no total computable procedure which, given $e$, returns a
finite list of objects meeting every isomorphism class of
$\mathcal G_e$. In particular there is no such primitive recursive
procedure, even though the presentations have decidable arrow validity
and a uniform bound of two on the number of isomorphism classes.
\end{theorem}

\begin{proof}
Suppose such a procedure existed, returning $L_e$. Add $0$ if necessary,
which preserves coverage, and compute the largest entry $H_e$.
Run the finite test $h(e,H_e)$. If it is one, the machine halts.
If it is zero, suppose instead that the machine halts at some later
time $t>H_e$. Every entry $a$ of $L_e$ satisfies $a\leq H_e$, so the
machine has not halted within $a$ steps and $h(e,a)=0$. But
$h(e,t)=1$. Thus $t$ is inequivalent to every member of $L_e$,
contradicting coverage. Therefore the zero answer implies that the
machine never halts.

The proposed procedure would decide halting. To recall the elementary
diagonal contradiction, a total decider for whether a program halts on
its own code would allow a program which loops when the decider says
that it halts and halts when the decider says that it loops. Applying
the decider to the code of that program contradicts either answer.
A blank-input decider yields this self-input decider by effectively
constructing a machine which first writes the fixed input and then
simulates the given program. Thus the assumed covering-list procedure
cannot exist.
\end{proof}

The theorem rules out a uniform algorithm; it is not an independence
claim about a particular elliptic curve or a particular presented
category. It also does not concern finite category tables, whose entire
object list is already available. The obstruction is precisely the
missing positive information in a bound on the number of components.

\subsection{A supplied finite cover gives an explicit skeleton}

Let $E$ be a fixed PR decidable equivalence relation on natural-number
objects, with its equivalence laws proved in iPRA. Suppose a finite list
$L$ is supplied together with the covering assertion
\begin{equation}
            \forall x\ \exists j<|L|\ E(x,L_j).
\label{eq:groupoid-cover}
\end{equation}
This assertion is an explicit hypothesis or a previously proved
theorem. It is not treated as a decidable property of the finite list.

\begin{theorem}[Finite-cover skeleton construction]
\label{thm:finite-cover-skeleton}
There is a PR algorithm producing a sublist $R$ of $L$ whose elements
are pairwise inequivalent and still satisfy \eqref{eq:groupoid-cover}.
The corresponding thin groupoid is explicitly equivalent to the
discrete finite category on the indices of $R$. The proof uses only
bounded scans and quantifier-free induction, followed by generalization
over the arbitrary object.
\end{theorem}

\begin{proof}
Scan $L$ in order. Retain the next entry exactly when it is inequivalent
to every entry retained so far. Since $E$ is decidable and the retained
list has length at most $|L|$, every test is bounded. Keep the following
finite invariant: the retained entries are pairwise inequivalent, and
each already scanned entry is equivalent to some retained entry.
The invariant holds for the empty prefix. A newly retained entry is
inequivalent to all its predecessors by the test. A discarded entry
has, by failure of the bounded test, an equivalent retained entry.
Thus the invariant is preserved. At the end it gives pairwise
inequivalence and coverage of the original list.

For an arbitrary object $x$, the hypothesis supplies an index $j$ with
$E(x,L_j)$. The finite invariant supplies an index $i$ with
$E(L_j,R_i)$. Transitivity gives $E(x,R_i)$. This proves coverage of
all objects without choosing from an unbounded family: define $q(x)$
by scanning the finite list $R$ for the first equivalent representative.
The covering proof shows that the scan succeeds. Formally totalize the
scan with a failure flag; its failure branch is ruled out by the same
argument. Since the object set contains $0$, coverage also proves
$|R|>0$.

Let $\mathcal D_R$ be the discrete category with object indices
$0,\ldots,|R|-1$. Define $F$ from the thin groupoid of $E$ to
$\mathcal D_R$ by $F(x)=q(x)$. If $E(x,y)$, then both representatives
$R_{q(x)}$ and $R_{q(y)}$ are equivalent to $x$, using symmetry and
transitivity. Pairwise inequivalence forces $q(x)=q(y)$, so $F$ maps
the arrow $x\to y$ to the identity of that common index. This also
proves functoriality directly.

Define $U:\mathcal D_R\to\mathcal G_E$ by $U(i)=R_i$ and by mapping
each identity to the corresponding identity. The first-match scan gives
$q(R_i)=i$, because no other retained entry is equivalent to $R_i$.
Consequently $FU$ is exactly the identity functor. For every object $x$,
the covering relation gives the unique arrow
\[
                       R_{q(x)}\longrightarrow x.
\]
Its inverse is the reversed endpoint pair. These arrows form a natural
isomorphism $UF\cong\mathrm{Id}$: both sides of every naturality
equation are arrows with the same endpoints, and the groupoid is thin.
This supplies the equivalence data explicitly.

The list invariant is a bounded Boolean formula. Its induction never
contains the unbounded covering assertion; that assertion remains a
hypothesis and is instantiated only after fixing $x$. Hence the
construction respects quantifier-free induction. When the thin groupoid
has a known finite object enumeration, that enumeration may itself be used as
$L$; no further choice principle is required.
\end{proof}

This positive theorem and Theorem~\ref{thm:no-two-class-skeleton}
locate the boundary sharply. A complete finite cover supplies enough
data for an explicit skeleton and equivalence. A numerical upper bound
on the number of components does not supply that cover. The same
discipline should be applied to finite lists, heights and exceptional
sets in the inspected arithmetic routes~\cite{repo-mw,repo-siegel,repo-faltings}.
